\documentclass[11pt]{article}
\usepackage[T1]{fontenc}
\usepackage[utf8]{inputenc}
\usepackage{lmodern,amsmath,amssymb,amsthm,mathtools}
\usepackage[margin=1in]{geometry}
\usepackage[hidelinks]{hyperref}
\newtheorem{theorem}{Theorem}
\newtheorem{lemma}[theorem]{Lemma}
\newtheorem{proposition}[theorem]{Proposition}
\newtheorem{corollary}[theorem]{Corollary}
\newcommand{\R}{\mathbb R}\newcommand{\Q}{\mathbb Q}
\newcommand{\Z}{\mathbb Z}\newcommand{\N}{\mathbb N}
\newcommand{\DR}{\mathrm{DR}}\newcommand{\ee}{\mathrm e}
\newcommand{\lie}{\mathfrak u}\newcommand{\id}{\mathrm{id}}
\newcommand{\HH}{\operatorname{Horiz}}\newcommand{\poly}{\operatorname{poly}}
\title{GTZ-LP-04: the petal quotient and degree--rank construction}
\author{Author proposal: /root/astra\_v14\_assertions\_arithmetic}
\date{7 November 2026}
\begin{document}\maketitle
\noindent\textbf{Status and scope.} This is a proposed project proof under
protocol 1.4, awaiting independent review. It proves the algebraic quotient,
its exact lattice and character statements, and the comparison needed by
\cite[Lemma 12.1]{GTZ}, conditional on the explicitly stated frequency input.
It does not prove that input or the entire inverse theorem. The erratum
\cite{Err} motivated the repair but is not a theorem premise. The related-use
audit at the end distinguishes completed connecting arguments from other
proof units. The rational-subgroup and tautification lemmas \cite{LP13,TAUT}
are separately versioned project proofs awaiting independent acceptance.

\section{Indices, free groups, and the actual objects}
Fix integers $m\geq1$, $1\leq R\leq m$, and finite disjoint sets of letters
$C_i,L_i,B_i$ for $1\leq i\leq m$. Their union is $\mathcal A$.
The letters in $C_i$ are core, those in $L_i$ retained linear petals, and
those in $B_i$ discarded petals. Every letter in these three sets has weight
$i$. A bracket word has length $\ell$ equal to its number of entries, weight
$w$ equal to the sum of their weights, and a multidegree recording the number
of occurrences of each individual letter. Repeated entries count repeatedly.
All bracketings are allowed, including length-one words. Write $p$ for the
number of petal entries, and $b$ for the number of discarded entries.
We use $[x,y]=x^{-1}y^{-1}xy$, real powers
$x^t=\exp(t\log x)$, and $\ee(t)=\exp(2\pi\mathrm i t)$.

Let $\mathcal F$ be the free real Lie algebra on symbols $X_a$, $a\in\mathcal A$.
Its usual length grading and its individual-letter multigrading are retained.
Take the homogeneous ideal spanned by all words for which
\begin{equation}\label{eq:truncate}
 w>m\quad\hbox{or}\quad(w=m\hbox{ and }\ell>R),
\end{equation}
and let $\lie$ be the quotient. These forbidden multidegrees are closed under
addition of any nonzero multidegree, so their span is indeed an ideal.
All surviving words have length at most $m$. Thus $\lie$ is finite-dimensional
and nilpotent. Let $U=\exp\lie$, with multiplication given by the finite BCH
polynomial, and set $e_a=\exp X_a$. The group is connected and simply connected.
These are the weighted universal objects intended in
\cite[Definition 9.1, p.~1279]{GTZ}, with the top rank in the relation fixed
to be $R$, not a running rank index.

Use $\DR=\{(d,r):d,r\in\N,\ 0\leq r\leq d\}$ with lexicographic order.
Put $U_{(0,0)}=U$. For $d>0$ let
\begin{equation}\label{eq:filtration}
 \lie_{(d,r)}=
 \operatorname{span}_{\R}\{W:w(W)>d\ \hbox{or}
 [w(W)=d,\ \ell(W)\geq\max(r,1)]\},
 \qquad U_{(d,r)}=\exp\lie_{(d,r)}.
\end{equation}
In particular the first alternative has \emph{no} length restriction.
For ranks larger than $d$ use the notation $U_{(d,r)}=U_{(d+1,0)}$ only
as the next-index convention; do not put such pairs into the index monoid.
Write $U_d=U_{(d,0)}$. The top group is
\[
 T=U_{(m,R)}=\exp\operatorname{span}_{\R}
       \{W:w(W)=m,\ \ell(W)=R\}.
\]

\begin{lemma}\label{lem:filtration}
The groups in \eqref{eq:filtration} form a taut degree--rank filtration of
degree at most $(m,R)$. Each is connected and normal in $U$, and $T$ is central.
\end{lemma}
\begin{proof}
Bracketing words adds weights, lengths, and individual-letter multidegrees.
Each displayed span is an ideal, since bracketing with a nonzero word increases
weight. It is a Lie subalgebra and BCH therefore identifies its exponential
with a connected subgroup. Nesting follows directly from lexicographic order:
passing to a larger degree removes all lower-weight words, and at fixed degree
passing to a larger rank increases the length cutoff.
For positive indices $(d,r),(e,t)$, a bracket of two allowed words either has
weight greater than $d+e$, or has weight $d+e$ and length at least
$\max(r,1)+\max(t,1)\geq\max(r+t,1)$. Hence the Lie bracket lies in
$\lie_{(d+e,r+t)}$. The BCH commutator formula has this bracket and further
brackets, all in the same ideal, which proves the group commutator condition.
At index zero use normality. Every generator has weight at least one, so
$U_{(1,0)}=U$; also $U_{(d,0)}=U_{(d,1)}$ for all $d>0$.
The truncation gives the claimed degree support. Bracketing a top word with
any generator raises its weight above $m$, proving centrality of $T$.
\end{proof}

\section{An integral coordinate lemma}
Here is the integral premise used, with its scope separated from the subsequent
quotient argument. Hall collection says that in a free nilpotent group on a
finite ordered alphabet, every element has a unique collected expression in
the basic commutators with integer exponents; the length-$j$ basic commutators
give a basis of its $j$th lower-central quotient. This published statement is
recorded in \cite[Definition 5.4.3, Theorem 5.4.5 and the final paragraph
of \S5.4, pp.~96--98]{Tointon}. The weight in that reference is bracket
\emph{length}, not the additional weight $w$ used here.

The corresponding basic Lie brackets form a homogeneous basis of the free
nilpotent Lie algebra. To spell out the connection, antisymmetry and Jacobi
reorder a bracket $[c_i,c_j]$ until $i>j$, and when
$c_i=[c_s,c_t]$ until $j\geq t$. Jacobi replaces the offending bracket by
brackets with the smaller inner entry moved outward. Induction on total
length and then the finite Hall ordering gives spanning by basic brackets.
Independence follows by mapping the free Lie algebra to the real tensor of
the lower-central graded Lie ring of the free group: the bracket is induced
by the group commutator, and Hall's basis assertion makes the images of those
basic brackets independent in each length. This also respects the individual
letter multigrading because the reordering identities preserve every count.

\begin{lemma}[Upward-closed monomial quotients]\label{lem:hall}
In a free nilpotent Lie algebra of class at most $m$, discard an upward-closed
set of nonzero individual-letter multidegrees. Let $V$ be the exponential
of the quotient and $\Lambda$ the discrete group generated by the surviving
$\exp X_a$. The surviving Hall commutators give unique real coordinates
\[
 (t_1,\ldots,t_v)\longmapsto c_1^{t_1}\cdots c_v^{t_v}
\]
on $V$, and $\Lambda$ is exactly the set with all $t_i\in\Z$.
It is a lattice. A further upward-closed set of multidegrees defines a
rational connected normal subgroup $V_S$. The quotient map $q$ between
two such monomial quotients obeys, for each such $S$,
\begin{equation}\label{eq:saturation}
 q(\Lambda\cap V_S)=q(\Lambda)\cap q(V_S).
\end{equation}
\end{lemma}
\begin{proof}
Order basic commutators by increasing length. For a basic group commutator
$c$, BCH gives $\log c=W_c+$ terms of strictly greater length, where $W_c$
is the corresponding basic Lie bracket. Each correction uses all the letters
of $c$ and possibly more occurrences of them. A discarded basic word
therefore vanishes in the quotient, including every BCH correction.
The surviving leading brackets are a basis, since the discarded Lie ideal
is the direct sum of precisely the discarded multigraded components.
Consequently the logarithm of the displayed coordinate map is triangular
by length with diagonal entries one. Solving successively in each length
gives a polynomial inverse and proves unique real coordinates.

Hall collection expresses every word in the original $\exp X_a$ and their
inverses as a product of basic group commutators with integer exponents.
After quotienting, discarded factors are identities, so it gives the claimed
integer coordinates. Conversely every surviving basic commutator is itself
a word in the generators, which proves the other inclusion. In particular
the subgroup is discrete. Suffix spans in the length-ordered basis are ideals,
since every nonzero bracket increases length. Right multiplication by
$c_i^{-n}$ changes coordinate $i$ by $-n$, leaves all earlier coordinates
unchanged, and can change only later coordinates. Reduce the coordinates
successively modulo integers. Every right coset has a representative with
coordinates in $[0,1]^v$, a compact set, proving cocompactness.

For $V_S$, quotient once more by its multigraded ideal. The same triangular
coordinate argument says that an element belongs to $V_S$ exactly when all
coordinates outside $S$ vanish. Thus its lattice consists of the integer
coordinates on the basic commutators in $S$, and the same compact-coordinate
argument proves rationality. A quotient deletes discarded basic factors and
keeps the remaining ones in their original order. Membership in $q(V_S)$
forces all retained coordinates outside $S$ to vanish. An element of
$q(\Lambda)\cap q(V_S)$ therefore lifts as the identical integer product
of retained basic commutators in $S$. This proves \eqref{eq:saturation}.
\end{proof}

The truncation \eqref{eq:truncate} is upward-closed, as are all the thresholds
\eqref{eq:filtration}. The lemma proves that
$\Gamma=\langle e_a:a\in\mathcal A\rangle_{\Z}$ is a lattice in $U$, and
every $U_{(d,r)}$ is rational. It also proves the lattice statement below
without treating it as a general property of arbitrary rational quotients.

\section{The corrected kernel and the descended character}
Call a word bad when $p\geq2$ or $b\geq1$. Let $I\subseteq\lie$ be the
span of all bad Lie words of every weight. Let $G^*$ be the connected real
subgroup generated by all bad iterated \emph{group} commutators in the $e_a$,
again of every weight and including length one. Equivalently, the real
generated subgroup contains every real power of each such commutator.
This convention is essential: a discrete generated subgroup would not be
the kernel of a map of simply connected Lie groups.

Let $\widetilde{\lie}$ be the free Lie algebra on $C_i\cup L_i$ with relations
\eqref{eq:truncate} and the vanishing of all words with at least two petal
entries. Give $\widetilde G=\exp\widetilde{\lie}$ the filtration
\eqref{eq:filtration} on its surviving words. Put
$\widetilde\Gamma=\langle\widetilde e_a:a\in C_i\cup L_i\rangle_{\Z}$.
Define $q:U\to\widetilde G$ by $e_a\mapsto\widetilde e_a$ on retained
letters and $e_a\mapsto1$ on discarded letters.

\begin{theorem}\label{thm:quotient}
The homomorphism $q$ is well-defined and surjective, and
\[
 \ker q=G^*=\exp I.
\]
Thus $G^*$ is connected, rational and normal; it need not be central. For
every degree--rank index $a$, including zero,
\begin{align}
 q(U_a)&=\widetilde G_a\simeq U_a/(G^*\cap U_a),\label{eq:filteredquot}\\
 q(\Gamma\cap U_a)&=\widetilde\Gamma\cap\widetilde G_a.\label{eq:exactlattice}
\end{align}
Suppose $\eta:T\to\R$ is a continuous homomorphism integral on
$T\cap\Gamma$, and for every bad group commutator $c$ of weight $m$ and
length $R$ one has $\eta(c)=0$. Then there is a unique continuous
$\widetilde\eta:\widetilde G_{(m,R)}\to\R$ satisfying
\[
 \widetilde\eta(q(t))=\eta(t)\quad(t\in T),\qquad
 \widetilde\eta(\widetilde\Gamma\cap\widetilde G_{(m,R)})\subseteq\Z.
\]
\end{theorem}
\begin{proof}
Bad multidegrees are upward-closed, so $I$ is an ideal. The universal property
of the free Lie algebra identifies $\lie/I$ with $\widetilde{\lie}$:
killing discarded letters kills every word involving one, and killing
all words with two petals gives exactly the other defining relations.
There are no additional relations because both presentations impose the
same homogeneous ideal on the same free Lie algebra. Exponentiation of this
surjective Lie map gives $q$, and its kernel is $\exp I$ by the global
exponential diffeomorphisms.

BCH puts the logarithm of every bad group commutator in $I$. Conversely
take the bad Hall commutators. Their logarithms have the bad basic Lie
brackets as leading terms, with only higher-length bad terms. The triangular
matrix on the finite-dimensional bad subspace is invertible. Their logarithms
therefore span $I$, and the connected real subgroup they generate is $\exp I$.
This proves the asserted description of $G^*$, not just its normal closure.
All these spaces are spanned by rational basic words. Lemma~\ref{lem:hall}
proves rationality and $q(\Gamma)=\widetilde\Gamma$.

The Lie map deletes bad multigraded components and preserves weight and
length of every surviving word. Applying \eqref{eq:filtration} before or
after the map therefore gives identical spaces. Exponentiating proves
\eqref{eq:filteredquot}, and Lemma~\ref{lem:hall} proves
\eqref{eq:exactlattice}.

The direct multigrading gives
\[
 I\cap\log T=\operatorname{span}_{\R}
    \{W:W\hbox{ bad},\ w(W)=m,\ \ell(W)=R\}.
\]
There is no cancellation between a bad multidegree and a good one. At
weight $m$ the BCH correction terms in a group commutator have larger
weight and vanish. Thus the bad group commutators of weight $m$ and
length $R$ are precisely the exponentials of their Lie brackets.
The assumption on $\eta$, and real linearity of a continuous homomorphism
on the vector group $T$, prove that $\eta$ vanishes on $G^*\cap T$.
It consequently descends through \eqref{eq:filteredquot} at the top index.
Finally \eqref{eq:exactlattice} lifts every top lattice point downstairs to
one upstairs, where $\eta$ is integral. This proves the last assertion.
\end{proof}

\paragraph{Why the omitted lower weights matter.}
With one core letter $x$ and one discarded letter $b$, both of weight one,
and $m=3,R=3$, the quotient kills $b$ itself. A subgroup generated only by
bad weight-three words cannot contain $b$. Moreover $[b,x]$ survives in $U$,
so the actual kernel is not central. This example tests both false shortcuts.

\paragraph{The precise analytic input.}
In the notation of \cite{GTZ}, take $m=s-1$, $R=r_*$,
$|C_i|=D_{*,i}$, $|L_i|=D'_{\mathrm{lin},i}$ and
$|B_i|=D'_{\mathrm{ind},i}$. Use the \emph{final output} of Lemma 10.10,
after rational-core and small frequencies have been removed by Lemma 10.9.
Then the two hypotheses on $\eta$ in Theorem~\ref{thm:quotient} are exactly
Theorem 11.1(i),(ii): all weight-$m$, length-$R$ words with two petals or
one beyond-linear petal are annihilated. The repeated (iii),(iv) add nothing.
For a bracketing convention restricted to left-normed words, Jacobi expresses
any other bracket of these same entries as a linear combination of such words;
their petal counts and total weight are unchanged. At the central top layer
this is also an exact identity of logarithms, so the vanishing extends to
all bracketings used here.

Theorem 11.1 requires the hypotheses of Theorem 7.2, the limit family of
degree--rank $(m,R)$ nilcharacters, and the sunflower representation with
its common standard dimension vector and vertical frequency. Its proof uses
the additive-quadruple bias conclusion of Proposition 7.3 and the stated
independence from Lemma 10.10. In the repaired project, that bias conclusion
must come from LP-03, not the false original Proposition 8.3. No use of
Theorem 1.3 is made to supply it. This unit assumes these inputs; merely
partitioning generators does not prove the required vanishing.

\section{Vertical test functions and petal actions}
\begin{lemma}\label{lem:vertical}
For the descended integral character there is a finite-dimensional
unit-vector-valued smooth function on $\widetilde G/\widetilde\Gamma$ of
vertical frequency $\widetilde\eta$. Its dimension and Lipschitz bound depend
only on the fixed filtered group, lattice, metric and character.
\end{lemma}
\begin{proof}
The central torus $\widetilde T/(\widetilde T\cap\widetilde\Gamma)$ acts
freely. Choose finitely many local smooth sections of the resulting principal
torus bundle and smooth real functions $a_j$ subordinate to their domains,
with $\sum_j a_j^2>0$ everywhere on the base. Normalize them so
$\sum_j a_j^2=1$. In a local trivialization $(z,t)$ put the $j$th component
equal to $a_j(z)\ee(\widetilde\eta(t))$, extending by zero outside its
support. The support stays inside the chart, so the extension is smooth.
Transition functions change only the component's phase. The vector has
norm one, and all components have the same vertical frequency.
Compactness gives a finite Lipschitz bound. If the torus or character is
trivial the same construction, or the constant function, applies.
\end{proof}

The quotient Lie algebra is a direct sum $\mathfrak c\oplus P$, where
$\mathfrak c$ consists of words with no petals and $P$ of words with one.
Here $P$ is an abelian ideal: a bracket of two such words has two petals
and is zero. Thus $\widetilde G_{\mathrm{petal}}=\exp P$ is abelian, normal
and rational by Lemma~\ref{lem:hall}. Decompose $P=\bigoplus_{a\in L}P_a$
by the unique petal letter. Core conjugation preserves each $P_a$.
For a tuple $t=(t_a)\in\R^L$, the map fixing core generators and sending
$X_a$ to $t_aX_a$ on linear petals is a Lie endomorphism $M_t$:
it preserves all defining ideals. It scales $P_a$ by $t_a$, so on the petal
group, written multiplicatively,
\[
 M_tM_u=M_{tu},\qquad M_t(p)M_u(p)=M_{t+u}(p),\qquad
 M_t(gpg^{-1})=gM_t(p)g^{-1}.
\]
The last identity follows because $M_t$ commutes with every core adjoint
action and conjugation by $P$ acts trivially on $P$. This proves the precise
identities used in (12.8)--(12.10), without treating $M_t$ as an automorphism
when some $t_a=0$. For integral tuples $t$, $M_t$ preserves
$\widetilde\Gamma$, because it sends each generator to an integer power.
It also preserves the petal lattice by its subgroup invariance. These are
the lattice facts needed by the subsequent shear semidirect products.

\section{The comparison filtration and corrected Lemma 12.1}
We record sufficient input explicitly. Let $G/\Lambda$, $G^0/\Lambda^0$
be fixed taut rational degree--rank nilmanifolds of degrees at most $(m,R)$
and strictly below $(m,R)$, respectively. Let $\phi:U\to G$ be filtered
and send $\Gamma$ into $\Lambda$. Suppose $a_h,c_h$ are polynomial sequences
and $F_h$ are unit-vector-valued uniformly Lipschitz functions on
$G/\Lambda\times G^0/\Lambda^0$ such that
\begin{equation}\label{eq:Fchar}
 F_h(\phi(t)x,z)=\ee(\eta(t))F_h(x,z)\quad(t\in T).
\end{equation}
Let $E_{h,i}=\prod_{a\in C_i\cup L_i\cup B_i}e_a^{\xi_{h,a}}$ in a fixed
order, and suppose the horizontal Taylor coefficient of $a_h\Lambda$ is
that of $\phi(E_{h,i})$ in $G_{(i,1)}/G_{(i,2)}$ modulo the horizontal
lattice. Let
\[
 b_h(n)=\prod_{i,a\in C_i}\widetilde e_a^{\xi_{h,a}\binom ni}
           \prod_{i,a\in L_i}\widetilde e_a^{\xi_{h,a}\binom ni}.
\]
These are exactly the core and retained-petal factors in (12.3)--(12.5).
Using the function from Lemma~\ref{lem:vertical}, set
\[
 \chi_h(n)=F_h(a_h(n)\Lambda,c_h(n)\Lambda^0),\qquad
 \widetilde\chi_h(n)=\widetilde F(b_h(n)\widetilde\Gamma).
\]

\begin{proposition}\label{prop:comparison}
Under these hypotheses, every component of
$\chi_h\otimes\overline{\widetilde\chi_h}$ is a nilsequence of
degree--rank strictly below $(m,R)$. The representing filtered nilmanifold
is fixed independently of $h$ and the coefficients; its Lipschitz bounds
depend only on the fixed data and original bounds.
\end{proposition}
\begin{proof}
For $d>0$ put $\bar r=\max(r,1)$ and let $\sigma(d,\bar r)$ denote the
successor of $(d,\bar r)$ in the positive degree--rank list. In the product
$G\times G^0\times\widetilde G$, let $H_{(d,r)}$ be the connected subgroup
generated by
\begin{align*}
 J_{(d,r)}&=G_{\sigma(d,\bar r)}\times G^0_{(d,\bar r)}
                        \times\widetilde G_{\sigma(d,\bar r)},\\
 D(U_{(d,\bar r)})&=\{(\phi(u),1,q(u)):u\in U_{(d,\bar r)}\}.
\end{align*}
Set $H_0=H_{(1,0)}$. This specifies the zero group; it is not an arbitrary
full product. All generator subgroups are rational. For the diagonal image,
the image of the appropriate lattice generates its real subgroup by the
Hall coordinates and $\phi(\Gamma)\subseteq\Lambda$. The rational-generated
subgroup lemma of \cite[Section 2]{LP13} applies.

Nesting is direct. For positive indices $i,j$, commutators of two diagonal
generators lie in $D(U_{i+j})$. A commutator involving a $J_i$ generator
has its first and third coordinates in the successor group at $i+j$,
since $\sigma(i)+j\succeq\sigma(i+j)$, with the convention replacing
rank zero by one. Its middle coordinate is in $G^0_{i+j}$. The same
argument works for two $J$ generators. Taking generated Lie algebras and
BCH gives $[H_i,H_j]\subseteq H_{i+j}$. At zero use
$H_0=H_{(1,0)}=H_{(1,1)}$ and nesting. Thus this is a taut rational
degree--rank filtration. Its top group is exactly
\begin{equation}\label{eq:Htop}
 H_{(m,R)}=D(T),
\end{equation}
because every group in $J_{(m,R)}$ is trivial.

We next check actual polynomial membership, including constant coefficients
and the distinction between a horizontal vector and its torus class.
Write $a_h(0)=v_h\lambda_h$ with $v_h$ in a fixed compact representative
set and $\lambda_h\in\Lambda$. Replace $a_h(n)$ by
$v_h^{-1}a_h(n)\lambda_h^{-1}$, and replace $F_h(x,z)$ by $F_h(v_hx,z)$.
Do the same in the second coordinate. Values of $\chi_h$ are unchanged,
the new sequences start at the identity, and Lipschitz bounds stay uniform.
Since $G$ is taut, conjugation by $v_h$ acts trivially on every
$G_{(i,1)}/G_{(i,2)}$: its commutator is in
$G_{(i+1,1)}\subseteq G_{(i,2)}$. The horizontal torus conditions remain
unchanged. Equation \eqref{eq:Fchar} remains unchanged since the top group
of the taut group is central. Tautify the original targets before constructing
the total-frequency representation, using \cite{TAUT} and then
\cite[Lemma 9.12]{GTZ}; no conjugated, potentially irrational replacement
of $\phi$ is being silently used here.

For each $i$, choose $\lambda_{h,i}\in\Lambda\cap G_i$ lifting the
horizontal lattice difference between the new $a_h$ and $\phi(E_{h,i})$.
Such a lift exists by the definition of the horizontal lattice. Multiply
$a_h(n)$ on the right by the inverse of
$\prod_i\lambda_{h,i}^{\binom ni}$. This does not change its orbit.
The factor is polynomial by the Taylor criterion
\cite[Lemma B.9]{GTZ}. The horizontal Taylor map is a homomorphism
\cite[Definition 9.6 and the argument following it]{GTZ}, so now the
horizontal equalities hold in the vector groups themselves.

Taylor-expand $a_h(n)=\prod_{i\geq1}a_{h,i}^{\binom ni}$. Then
$a_{h,i}\phi(E_{h,i})^{-1}\in G_{(i,2)}$. Hence
$(a_{h,i},1,q(E_{h,i}))\in H_{(i,1)}$. The Taylor criterion and product
closure \cite[Lemma B.9 and Corollary B.4]{GTZ} show that
\[
 n\mapsto\left(a_h(n),1,\prod_iq(E_{h,i})^{\binom ni}\right)
\]
is polynomial for $H$. Both its last coordinate and $b_h$ have the same
horizontal Taylor coefficients, since multiplication of factors is addition
in each horizontal quotient and discarded generators have image one.
Their ratio is polynomial and all its horizontal Taylor coefficients vanish.
Its Taylor coefficient at index $i$ therefore lies in
$\widetilde G_{(i,2)}$. These groups, with zeroth group
$\widetilde G_{(1,2)}$, form a degree filtration by the same commutator
check. The ratio starts at one; Lemma B.9 now shows that it is polynomial
for this smaller filtration. Embedded in the third coordinate it is
polynomial in $H$. The second-coordinate map $(1,c_h(n),1)$ is also
polynomial in $H$. Product closure proves that the exact triple
$(a_h(n),c_h(n),b_h(n))$ is polynomial in $H$.

Restrict $F_h\otimes\overline{\widetilde F}$ to
$H/(H\cap(\Lambda\times\Lambda^0\times\widetilde\Gamma))$.
On $D(t)$ its vertical multiplier is
$\ee(\eta(t)-\widetilde\eta(q(t)))=1$ for every $t\in T$.
It therefore descends through the central rational subgroup
\eqref{eq:Htop}. The central quotient and lattice argument in
\cite[Sections 2 and 6]{LP13} applies here: this kernel is contained in
every nontrivial indexed subgroup, unlike the earlier kernel $G^*$.
The top group becomes trivial, so the quotient has degree--rank strictly
below $(m,R)$. Restriction, quotient charts, and the compact translations
cost only constants from the fixed geometry. This proves the proposition.
\end{proof}

When $R=1$, strictly below $(m,1)$ means ordinary degree at most $m-1$.
For $R>1$ it means degree--rank at most $(m,R-1)$, not necessarily degree
at most $m-1$. The first sentence of the published proof of Lemma 12.1
claims the latter unnecessarily strong conclusion. The proved former
conclusion is exactly what the following rank-induction step consumes.

\section{Uniformity, limits, and the remaining Section 12 interfaces}
Every group, ideal, Hall list, lattice and quotient above depends only on
$m,R$ and the fixed finite labelled dimension data. All formulas are finite
polynomials, Lie-group maps, or choices from fixed lattice and compact
representative sections. A fixed ordering chooses the first lift of a lattice
class; a fixed finite atlas chooses a bounded representative. Applying these
standard choices to representative sequences gives limit maps. The derivative
and vertical identities are required on their full ultrapower domains,
as in \cite{TAUT}; identities only at standard points are insufficient.
Thus no bound depends on $N$, on $h$, or on the real polynomial coefficients.
The character is fixed before choosing a vertical test function. No assertion
of a bound uniform over unbounded characters is made. Taking the dense
subfamily with fixed geometric data and character is the separate selection
step in Lemma 10.10 and Section 12, using Appendix A.

The ordinary degree filtration used to build the final semidirect product
is $\widetilde G_d=\widetilde G_{(d,0)}$. The petal ideal and scalar-action
identities proved above give precisely the normality, commutativity, and
integral stability used on pp.~1312--1314. In particular a bracket with two
petal entries is zero at \emph{every} weight. This assertion would fail if
one quotiented only by the old top-weight kernel.

There is a separate unit-norm issue at the cutoff in (12.11) and the paragraph
following it. A compactly supported scalar cutoff times a unit-vector function
is not itself unit-vector-valued. Here is an exact completion for that use.
On the final compact nilmanifold choose a smooth vertical-equivariant local
unit vector $V$ on the torus-coordinate chart, agreeing with
$\widetilde F(g\widetilde\Gamma)$ there. Choose a smooth function
$\theta$ of the torus coordinates equal to zero on the closed inner box and
equal to $\pi/2$ outside a slightly larger box still inside the chart.
These coordinates are unchanged by the top vertical action. Construct a
global unit-vector function $V_0$ of the same vertical character by
Lemma~\ref{lem:vertical}, applied to the final top torus. Then
\[
 \bigl(\cos\theta\,V,\ \sin\theta\,V_0\bigr)
\]
is smooth, unit-vector-valued, has the required vertical character, and equals
$(V,0)$ on the inner box. The first component extends by zero, since its
support stays inside the chart. Choose $\theta$ constant near both boundaries,
so both extensions are smooth. This preserves the required values for the
dense parameter set after shrinking to the inner box. Its geometry and bounds
are fixed before $N$. The final semidirect-product filtration and polynomial
orbit must still be checked at their displayed definitions; they are not
consequences of the character descent alone.

\section{Bounded audit of affected uses}
The failure mechanism is restriction of a generated ideal or filtration to
too few weights, followed by an unjustified quotient or lattice identification.
The audit scope is the universal degree--rank definitions and the consumers
of their commutator/Taylor quotients in Sections 9--12. It includes the
zero-index and repeated normalization assertions in those passages.
\begin{enumerate}
\item Definition 9.1, p.~1279: use the fixed $R=r_*$ in the top relation,
and count the $\ell$ actual entries. Equations \eqref{eq:truncate} and
\eqref{eq:filtration}, including all higher weights, give the intended
object. Lemmas~\ref{lem:filtration} and \ref{lem:hall} prove its filtration
and lattice, rather than inheriting typographical indices.
\item Definition 9.6, pp.~1281--1282: use
$\HH_i=G_{(i,1)}/G_{(i,2)}$, and $G_{(1,2)}=G_{(2,0)}$.
The vector/torus distinction is discharged in
Proposition~\ref{prop:comparison}, not by replacing a congruence by equality.
\item Definition 9.11 and Lemma 9.12, pp.~1284--1287: their $G_i$ means
$G_{(i,0)}$. First tautify, then choose lattice generators and construct
the filtered map. This order preserves a genuine lattice map; the proof
does not assume arbitrary real conjugation preserves the lattice.
\item Lemma 9.14 and Corollary 9.15, pp.~1287--1288, are explicitly stated
there not to be used later. Their comparison-filtration argument has the
same zero-index and horizontal-lattice issues. The construction in
Proposition~\ref{prop:comparison} supplies that mechanism with $q$ trivial
when the character is zero. Their broader reuse requires this version,
not the unqualified original paragraph.
\item Lemma 10.8, p.~1297: the vertical domain is
$U_{(m,R)}$, not the larger ordinary degree group $U_m$.
This unit consistently defines vertical frequencies only on that top group.
The source argument's other linear-algebra steps are not newly reproved here.
\item Lemma 10.10 and Theorem 11.1, pp.~1299--1307: use the final
core/linear/independent decomposition and the exact two vanishing assertions.
The bias and independence premises remain named upstream inputs.
The repeated (iii),(iv) of Theorem 11.1 are redundant. Reordering linear
and independent petals requires relabelling every associated index, not
changing which generators are annihilated.
\item Quotient construction, p.~1309: Theorem~\ref{thm:quotient} proves
the full kernel, every filtered quotient, and the top-lattice saturation
needed for integral vertical descent. The kernel is not asserted central.
\item Lemma 12.1, pp.~1309--1311: Proposition~\ref{prop:comparison}
proves the rank drop with actual constant normalization, horizontal lattice
lifts, exact polynomial orbit and conjugate vertical cancellation.
The degree-only strengthening is not used or claimed.
\item Petal actions, pp.~1311--1314: the proved abelian ideal and its
individual-petal splitting verify (12.8)--(12.10), including zero and
negative scaling parameters and integral lattice stability. The displayed
final multidegree filtration, chart gluing, and polynomial orbit remain
explicit downstream construction checks for LP-00/LP-07. The unit-norm
cutoff completion above removes the identified local norm obstruction.
\end{enumerate}
This list is an authored bounded audit, not independent source-first coverage
or acceptance of Sections 9--12. It does not classify unrelated errors outside
this mechanism. The separate project proofs and the assembled inverse theorem
must be independently reviewed at their exact versions.

\begin{thebibliography}{9}
\bibitem{GTZ} B. Green, T. Tao and T. Ziegler,
An inverse theorem for the Gowers $U^{s+1}[N]$-norm,
Annals of Mathematics 176 (2012), 1231--1372.
\url{https://doi.org/10.4007/annals.2012.176.2.11}.
\bibitem{Tointon} M. C. H. Tointon,
Introduction to Approximate Groups, LMS Student Texts 94,
Cambridge University Press, 2020, \S5.4, pp.~96--98.
\url{https://doi.org/10.1017/9781108652865}.
The author-posted published-book PDF is pinned in the source manifest;
its author page lists corrections, none affecting this interface.
\bibitem{LP13} Project proof by /root/astra\_v14\_baseline\_theta,
Lemma 13.2 repair, version lemma13-2-002, 7 November 2026.
Exact manifest and proof hashes are in the accompanying source-use record.
\bibitem{TAUT} Same author, Tautification and the shifted nilcharacter
construction, version tautification-003, 7 November 2026.
\bibitem{Err} B. Green, T. Tao and T. Ziegler,
Erratum for the inverse theorem, six-page author document,
pp.~5--6. Attribution and problem identification only; no correction theorem
is imported from this unpublished document.
\end{thebibliography}
\end{document}
